% !TeX root = ./forallxyyc.tex
% Chapter on modal logic. Original author: Rob Trueman, York
% University, http://www.rtrueman.com/

\part{Lógica modal}
\label{ch.ML}
\addtocontents{toc}{\protect\mbox{}\protect\hrulefill\par}

\chapter{Introdução à lógica modal}
\label{Intro}

A lógica modal (LM) é a lógica das \emph{modalidades}, maneiras pelas quais uma afirmação pode ser verdadeira. \emph{Necessidade} e \emph{possibilidade} são duas dessas modalidades: uma afirmação pode ser verdadeira, mas também pode ser necessariamente verdadeira (verdadeira independentemente de como o mundo poderia ter sido). Por exemplo, verdades lógicas não são apenas verdadeiras por causa de alguma característica acidental do mundo, mas verdadeiras aconteça o que acontecer. Uma afirmação possível pode não ser de fato verdadeira, mas poderia ter sido verdadeira. Usamos $\ebox$ para expressar necessidade e $\ediamond$ para expressar possibilidade. Assim, você pode ler $\ebox \metav{A}$ como \emph{É necessariamente o caso que} $\metav{A}$, e $\ediamond \metav{A}$ como \emph{É possivelmente o caso que} $\metav{A}$.

Há muitos tipos diferentes de necessidade. É \emph{humanamente impossível} que eu corra a 160 km/h. Dado o tipo de criaturas que somos, nenhum ser humano consegue fazer isso. Mas ainda assim não é \emph{fisicamente impossível} que eu corra tão rápido. Ainda não temos a tecnologia, mas é certamente fisicamente possível trocar minhas pernas biológicas por pernas robóticas capazes de correr a 160 km/h. Em contraste, é fisicamente impossível que eu corra mais rápido que a luz. As leis da física impedem que qualquer objeto acelere até essa velocidade. Mas nem isso é \emph{logicamente} impossível. Não é uma contradição imaginar que as leis da física pudessem ter sido diferentes e permitido que os objetos se movessem mais rápido que a luz.

Com que tipo de modalidade a LM lida? \emph{Com todas elas!} A LM é uma ferramenta muito flexível. Começamos com um conjunto básico de regras que governam $\ebox$ e $\ediamond$ e depois acrescentamos mais regras para adequá-las ao tipo de modalidade que nos interessa. Na verdade, a LM é tão flexível que nem precisamos pensar em $\ebox$ e $\ediamond$ como expressando \emph{necessidade} e \emph{possibilidade}. Poderíamos ler $\ebox$ como expressando \emph{demonstrabilidade}, de modo que $\ebox\metav{A}$ significa \emph{É demonstrável que} $\metav{A}$, e $\ediamond\metav{A}$ significa \emph{Não é refutável que} $\metav{A}$. De modo semelhante, podemos interpretar $\ebox$ como significando que $S$ \emph{sabe que $\metav{A}$} ou que $S$ \emph{acredita que $\metav{A}$}. Ou poderíamos ler $\ebox$ como expressando \emph{obrigação moral}, de modo que $\ebox \metav{A}$ significa \emph{É moralmente obrigatório que} $\metav{A}$, e $\ediamond \metav{A}$ significa \emph{É moralmente permitido que} $\metav{A}$. Tudo o que precisaríamos fazer seria elaborar as regras corretas para essas diferentes leituras de $\ebox$ e $\ediamond$.

Uma fórmula modal é aquela que inclui operadores modais como $\ebox$ e $\ediamond$. Dependendo da interpretação que atribuímos a $\ebox$ e $\ediamond$, diferentes fórmulas modais serão demonstráveis ou válidas. Por exemplo, $\ebox \metav{A} \eif \metav{A}$ pode dizer que \textquotedblleft se $\metav{A}$ é necessário, então é verdadeiro\textquotedblright{}, se $\ebox$ for interpretado como necessidade. Pode expressar \textquotedblleft se $\metav{A}$ é conhecido, então é verdadeiro\textquotedblright{}, se $\ebox$ expressar verdade conhecida. Sob ambas as interpretações, $\ebox\metav{A} \eif \metav{A}$ é válida: todas as proposições necessárias são verdadeiras aconteça o que acontecer e, portanto, verdadeiras no mundo atual. E, se uma proposição é conhecida como verdadeira, ela tem de ser verdadeira (não se pode conhecer algo falso). No entanto, quando $\ebox$ é interpretado como \textquotedblleft acredita-se que\textquotedblright{} ou \textquotedblleft deve ser o caso que\textquotedblright{}, $\ebox\metav{A} \eif \metav{A}$ não é válida: podemos acreditar em proposições falsas. Nem toda proposição que deveria ser verdadeira é de fato verdadeira; por exemplo, \textquotedblleft todo assassino será levado à justiça\textquotedblright{}. Isso \emph{deveria} ser verdadeiro, mas não é.

Consideraremos diferentes tipos de sistemas de LM. Eles diferem nas regras de prova permitidas e na semântica que usamos para definir nossas noções lógicas. Os diferentes sistemas que consideraremos são chamados de \mlK, \mlT, \mlSfour{} e \mlSfive. \mlK{} é o sistema básico; tudo o que é válido ou demonstrável em \mlK{} também é demonstrável nos outros. Mas há coisas que \mlK{} não prova, como a fórmula $\ebox A \eif A$ para a letra sentencial~$A$. Portanto, \mlK{} não é uma lógica modal apropriada para necessidade e possibilidade (em que $\ebox\metav{A} \eif \metav{A}$ deveria ser demonstrável). Isso é demonstrável no sistema \mlT, de modo que \mlT{} é mais apropriado para lidar com necessidade e possibilidade, mas menos apropriado para lidar com crença ou obrigação, pois nesses casos $\ebox\metav{A} \eif \metav{A}$ não deveria ser (sempre) demonstrável. Talvez o melhor sistema de LM para necessidade e possibilidade — e, em qualquer caso, o mais amplamente aceito — seja o mais forte dos sistemas que consideramos, \mlSfive.

\section{A linguagem da LM}
\label{TFLtoML}

Para fazer lógica modal, precisamos fazer duas coisas. Primeiro, queremos aprender a provar coisas na LM. Segundo, queremos ver como construir interpretações para a LM. Mas, antes de fazermos qualquer uma dessas coisas, precisamos explicar como construir sentenças na LM.

A linguagem da LM é uma extensão da LVF. Poderíamos ter começado com a LPO, o que nos daria a lógica modal quantificada (LMQ). A LMQ é muito mais poderosa que a LM, mas também é muito, muito mais complicada. Portanto, vamos manter as coisas simples e começar com a LVF.

Assim como a LVF, a LM começa com um estoque infinito de \emph{átomos}. Eles são escritos como letras maiúsculas, com ou sem subscritos numéricos: $A$, $B$, \dots, $A_1$, $B_1$, \dots. Em seguida, tomamos todas as regras sobre como formar sentenças da LVF e acrescentamos duas para $\ebox$ e~$\ediamond$:
\factoidbox{
\begin{factoidnumber}
\item Todo átomo da LM é uma sentença da LM.
\item Se $\metav{A}$ é uma sentença da LM, então $\enot\metav{A}$ é uma sentença da LM.
\item Se $\metav{A}$ e $\metav{B}$ são sentenças da LM, então $(\metav{A}\eand\metav{B})$ é uma sentença da LM.
\item Se $\metav{A}$ e $\metav{B}$ são sentenças da LM, então $(\metav{A}\eor\metav{B})$ é uma sentença da LM.
\item Se $\metav{A}$ e $\metav{B}$ são sentenças da LM, então $(\metav{A}\eif\metav{B})$ é uma sentença da LM.
\item Se $\metav{A}$ e $\metav{B}$ são sentenças da LM, então $(\metav{A}\eiff\metav{B})$ é uma sentença da LM.
\item Se $\metav{A}$ é uma sentença da LM, então $\ebox\metav{A}$ é uma sentença da LM.
\item Se $\metav{A}$ é uma sentença da LM, então $\ediamond\metav{A}$ é uma sentença da LM.
\item Nada mais é uma sentença da LM.
\end{factoidnumber}}
Eis alguns exemplos de sentenças da LM:
\begin{align*}
	& A\\
	& P\eor Q\\
	& \ebox A\\
	& C\eor \ebox D\\
	& \ebox\ebox (A\eif R)\\
	& \ebox\ediamond (S\eand (Z\eiff (\ebox W \eor \ediamond Q)))
\end{align*}

\chapter{Dedução natural para a LM}
\label{Proof}

Agora que sabemos formar sentenças na LM, podemos ver como \emph{provar} coisas na LM. Usaremos $\proves$ para expressar demonstrabilidade. Assim, $\metav{A}_1,\metav{A}_2, \dots \metav{A}_n \proves \metav{C}$ significa que $\metav{C}$ pode ser provada a partir de $\metav{A}_1,\metav{A}_2, \dots \metav{A}_n$. No entanto, examinaremos vários sistemas diferentes de LM e, por isso, será útil acrescentar um subscrito para indicar em qual sistema estamos trabalhando. Por exemplo, se quisermos dizer que podemos provar $\metav{C}$ a partir de $\metav{A}_1,\metav{A}_2, \dots \metav{A}_n$ \emph{no sistema}~\mlK, escreveremos: $\metav{A}_1,\metav{A}_2, \dots \metav{A}_n \proves_\mlK \metav{C}$.

\section{Sistema \mlK}
\label{K}

Começamos com um sistema particularmente simples chamado \mlK, em homenagem ao filósofo e lógico Saul Kripke. \mlK{} inclui todas as regras de dedução natural da LVF, tanto as regras derivadas quanto as básicas. Em seguida, \mlK{} acrescenta um tipo especial de subprova e duas novas regras básicas para $\ebox$.

Esse tipo especial de subprova se parece com uma subprova comum, exceto por ter um $\ebox{}$ em sua linha de hipótese em vez de uma fórmula. Nós as chamamos de \emph{subprovas estritas} — elas nos permitem raciocinar e provar coisas sobre possibilidades alternativas. O que podemos provar dentro de uma subprova estrita vale em qualquer possibilidade alternativa, em particular, em possibilidades alternativas nas quais as hipóteses vigentes em nossa prova podem não valer. Em uma subprova estrita, todas as hipóteses são desconsideradas e não podemos recorrer a nenhuma linha fora da subprova estrita (exceto quando permitido pelas regras modais dadas abaixo).

A regra $\ebox$I permite derivar uma fórmula $\ebox \metav{A}$ se pudermos derivar $\metav{A}$ dentro de uma subprova estrita. Ela é nosso método fundamental para introduzir $\ebox{}$ nas provas. A ideia básica é simples: se $\metav{A}$ é um teorema, então $\ebox \metav{A}$ também deve ser um teorema. (Lembre-se de que chamar $\metav{A}$ de teorema significa dizer que podemos provar $\metav{A}$ sem depender de hipóteses não descarregadas.)

Suponha que quiséssemos provar $\ebox(A\eif A)$. A primeira coisa que precisamos fazer é provar que $A\eif A$ é um teorema. Você já sabe como fazer isso usando a LVF. Basta apresentar uma prova de $A\eif A$ que não comece com premissas, como esta:
\begin{fitchproof}
		\open
		\hypo{1}{A}\AS
		\have{2}{A}\by{R}{1}
		\close
		\have{3}{A\eif A}\ci{1-2}
\end{fitchproof}
Mas, para aplicar $\ebox$I, precisamos ter provado a fórmula dentro de uma subprova estrita. Como nossa prova de $A \eif A$ não usa hipótese alguma, isso é possível.
\begin{fitchproof}
		\open
		\hypo{1}{\ebox}\AS
		\open
		\hypo{2}{A}\AS
		\have{3}{A}\by{R}{2}
		\close
		\have{4}{A\eif A}\ci{2-3}
		\close
		\have{5}{\ebox(A\eif A)}\boxi{1-4}
\end{fitchproof}

\factoidbox{
	\begin{fitchproof}
			\open
			\hypo[m]{m}{\ebox}\AS
			\have[n]{n}{\metav{A}}
			\close
			\have[\,]{o}{\ebox\metav{A}}\boxi{m-n}
	\end{fitchproof}
Nenhuma linha acima da linha $m$ pode ser citada por qualquer regra dentro da subprova estrita iniciada na linha $m$, a menos que a regra permita explicitamente essa citação.
}
É essencial enfatizar que, em uma subprova estrita, você não pode usar nenhuma regra que recorra a algo provado fora da subprova estrita. Há exceções, como a regra $\ebox$E abaixo. Essas regras declararão explicitamente que podem ser usadas dentro de subprovas estritas e citarão linhas fora da subprova estrita. Essa restrição é essencial; caso contrário, obteríamos resultados terríveis. Por exemplo, poderíamos oferecer a seguinte prova para justificar $A\therefore \ebox A:
\begin{fitchproof}
		\hypo{1}{A}\PR
		\open
		\hypo{2}{\ebox}\AS
		\have{3}{A}\by{uso incorreto de R}{1}
		\close
		\have{4}{\ebox A}\boxi{2-3}
\end{fitchproof}
Isso não é uma prova legítima, pois na linha 3 recorremos à linha 1, embora a linha 1 venha antes do início da subprova estrita na linha 2.

Dissemos acima que uma subprova estrita nos permite raciocinar sobre situações possíveis alternativas arbitrárias. O que pode ser provado em uma subprova estrita vale em todas as situações alternativas e, portanto, é necessário. Essa é a ideia por trás da regra $\ebox$I. Por outro lado, se supusemos que algo é necessário, assumimos com isso que é verdadeiro em todas as situações possíveis alternativas. Assim, temos a regra $\ebox$E:

\factoidbox{
	\begin{fitchproof}
			\have[m]{m}{\ebox\metav{A}}
			\open
			\hypo[\ ]{o}{\ebox}\AS
			\have[n]{n}{\metav{A}}\boxe{m}
			\close
		\end{fitchproof}
	\ebox E só pode ser aplicada se a linha $m$ (que contém $\ebox A$) estiver \emph{fora} da subprova estrita em que a linha $n$ ocorre, e essa subprova estrita não fizer parte de outra subprova estrita que não contenha~$m$.
		}
A regra $\ebox$E permite afirmar $\metav{A}$ dentro de uma subprova estrita se você tiver $\ebox \metav{A}$ fora dela. A restrição significa que você só pode fazer isso na primeira subprova estrita; não pode aplicar a regra $\ebox$E dentro de uma subprova estrita aninhada. Portanto, o seguinte não é permitido:
\begin{fitchproof}
	\have{1}{\ebox\metav{A}}
	\open
	\hypo{2}{\ebox}\AS
	\open
	\hypo{3}{\ebox}\AS
	\have{4}{\metav{A}}\by{uso incorreto de \ebox E}{1}
\close\close
\end{fitchproof}
O uso incorreto de $\ebox$E na linha~$4$ viola a condição, pois, embora a linha~$1$ esteja fora da subprova estrita em que a linha~$4$ ocorre, a subprova estrita que contém a linha~$4$ está dentro da subprova estrita iniciada na linha~$2$, que não contém a linha~$1$.

Comecemos com um exemplo.
\begin{fitchproof}
		\hypo{1}{\ebox A}\PR
		\hypo{2}{\ebox B}\PR
		\open
		\hypo{3}{\ebox}\AS
		\have{4}{A}\boxe{1}
		\have{5}{B}\boxe{2}
		\have{6}{A \eand B}\ai{4,5}
		\close
		\have{6}{\ebox(A \land B)}\boxi{3-6}
\end{fitchproof}
Também podemos misturar subprovas comuns e subprovas estritas:
\begin{fitchproof}
		\hypo{1}{\ebox (A \eif B)}\PR
		\open
		\hypo{2}{\ebox A}\AS
		\open
		\hypo{3}{\ebox}\AS
		\have{4}{A}\boxe{2}
		\have{5}{A \eif B}\boxe{1}
		\have{6}{B}\ce{4,5}
		\close
		\have{7}{\ebox B}\boxi{3-6}
		\close
		\have{8}{\ebox A \eif \ebox B} \ci{2-7}
\end{fitchproof}
Isso é chamado de \emph{regra de distribuição}, pois nos diz que $\ebox$ \textquotedblleft se distribui\textquotedblright{} sobre $\eif$.

As regras $\ebox$I e $\ebox$E parecem simples o suficiente — e, de fato, \mlK{} é um sistema muito simples! Mas \mlK{} é mais poderoso do que você talvez tenha imaginado. É possível provar várias coisas nele.

\section{Possibilidade}
\label{possibility}

Na última subseção, examinamos todas as regras básicas para \mlK. Mas talvez você tenha notado que todas essas regras tratavam de necessidade, $\ebox$, e nenhuma tratava de possibilidade, $\ediamond$. Isso ocorre porque podemos \emph{definir} possibilidade em termos de necessidade:
\factoidbox{
$\ediamond\metav{A}=_{df} \enot \ebox\enot \metav{A}$
}
Em outras palavras, dizer que $\metav{A}$ é \emph{possivelmente verdadeira} é dizer que $\metav{A}$ \emph{não é necessariamente falsa}. Como resultado, não é realmente essencial acrescentar $\ediamond$, um símbolo especial para possibilidade, ao sistema \mlK. Ainda assim, o sistema será \emph{muito} mais fácil de usar se o fizermos e, por isso, acrescentaremos as seguintes regras definicionais:
\factoidbox{
	\begin{fitchproof}
			\have[m]{m}{\enot\ebox\enot \metav{A}}
			\have[\, ]{n}{\ediamond \metav{A}}\diadf{m}
	\end{fitchproof}
	\begin{fitchproof}
			\have[m]{m}{\ediamond \metav{A}}
			\have[\, ]{n}{\enot\ebox\enot \metav{A}}\diadf{m}
	\end{fitchproof}
}
É importante não pensar nessas regras como uma adição real a \mlK: elas apenas registram a maneira como $\ediamond$ é definido em termos de $\ebox$.

Se quiséssemos, poderíamos encerrar aqui nossas regras para \mlK{}. Mas será útil acrescentar algumas regras de \emph{conversão modal}, que nos dão outras maneiras de alternar entre $\ebox$ e $\ediamond$:
\factoidbox{
	\begin{fitchproof}
			\have[m]{m}{\enot\ebox \metav{A}}
			\have[\, ]{n}{\ediamond \enot\metav{A}}\mc{m}
	\end{fitchproof}
	\begin{fitchproof}
			\have[m]{m}{\ediamond \enot \metav{A}}
			\have[\, ]{n}{\enot\ebox \metav{A}}\mc{m}
	\end{fitchproof}
	\begin{fitchproof}
			\have[m]{m}{\enot\ediamond \metav{A}}
			\have[\, ]{n}{\ebox \enot\metav{A}}\mc{m}
	\end{fitchproof}
	\begin{fitchproof}
			\have[m]{m}{\ebox\enot \metav{A}}
			\have[\, ]{n}{\enot\ediamond\metav{A}}\mc{m}
	\end{fitchproof}
}
Essas regras de conversão modal também não acrescentam poder ao \mlK, pois podem ser derivadas das regras básicas juntamente com a definição de $\ediamond$.

No sistema \mlK, usando Def\ediamond{} (ou as regras de conversão modal), podemos provar $\ediamond A \eiff \enot\ebox\enot A$. Ao apresentar o sistema \mlK, começamos com $\ebox$ como nosso símbolo modal primitivo e definimos $\ediamond$ em termos dele. Mas, se preferíssemos, poderíamos começar com $\ediamond$ como primitivo e então definir $\ebox$ da seguinte maneira: $\ebox\metav{A} =_{df} \enot \ediamond \enot \metav{A}$. Não há, portanto, nenhum sentido em que a necessidade seja de algum modo mais \emph{fundamental} que a possibilidade. Necessidade e possibilidade são igualmente fundamentais.

\section{Sistema \mlT}
\label{T}

Até aqui, concentramos nossa atenção em \mlK, um sistema modal muito simples. \mlK{} é tão fraco que nem sequer permite provar $\metav{A}$ a partir de $\ebox\metav{A}$. Mas, se estamos pensando em $\ebox$ como expressando \emph{necessidade}, queremos poder fazer esta inferência: se $\metav{A}$ é \emph{necessariamente verdadeira}, então certamente deve ser \emph{verdadeira}!

Isso nos leva a um novo sistema, \mlT, que obtemos acrescentando a seguinte regra a \mlK:
\factoidbox{
	\begin{fitchproof}
			\have[m]{m}{\ebox \metav{A}}
			\have[n]{n}{\metav{A}}\rt{m}
	\end{fitchproof}
	A linha $n$ na qual a regra R\mlT{} é aplicada não pode estar \emph{em} uma subprova estrita que comece depois da linha~$m$.
}

A restrição da regra~\mlT{} é, de certo modo, o oposto da restrição da regra $\ebox$E: você só pode usar $\ebox$E em uma subprova estrita aninhada, mas não pode usar \mlT{} em uma subprova estrita aninhada.

Podemos provar em \mlT{} coisas que não podíamos provar em \mlK, por exemplo, $\ebox A \eif A$.

\section{Sistema \mlSfour}
\label{S4}

\mlT{} permite eliminar as caixas de necessidade: de $\ebox \metav{A}$, você pode inferir $\metav{A}$. Mas e se quiséssemos acrescentar caixas extras? Isto é, passar de $\ebox\metav{A}$ para $\ebox\ebox\metav{A}$? Isso não seria problema se tivéssemos provado $\ebox\metav{A}$ aplicando $\ebox$I a uma subprova estrita de $\metav{A}$ que, por sua vez, não usasse $\ebox$E. Nesse caso, $\metav{A}$ é uma tautologia e, aninhando a subprova estrita dentro de outra subprova estrita e aplicando $\ebox$I novamente, podemos provar $\ebox\ebox \metav{A}$. Por exemplo, poderíamos provar $\ebox\ebox (P\eif P)$ assim:
\begin{fitchproof}
		\open
		\hypo{1}{\ebox}\AS
		\open
		\hypo{2}{\ebox}\AS
		\open
		\hypo{3}{P}\AS
		\have{4}{P}\by{R}{3}
		\close
		\have{5}{P\eif P}\ci{3-4}
		\close
		\have{6}{\ebox(P\eif P)}\boxi{2-5}
		\close
		\have{7}{\ebox\ebox(P\eif P)}\boxi{1-6}
\end{fitchproof}
Mas e se não tivéssemos provado $\ebox\metav{A}$ dessa maneira restrita, mas usado $\ebox$E dentro da subprova estrita de $\metav{A}$? Se colocarmos essa subprova estrita dentro de outra subprova estrita, a exigência da regra $\ebox$E de não citar uma linha que contenha $\ebox\metav{A}$ e esteja em outra subprova estrita ainda não concluída será violada. Ou, e se $\ebox\metav{A}$ fosse apenas uma hipótese com que começamos nossa prova? Poderíamos então inferir $\ebox\ebox\metav{A}$? Não em \mlT. E isso pode parecer uma limitação de \mlT{}, pelo menos se estivermos lendo $\ebox$ como expressando \emph{necessidade}. Parece intuitivo que, se $\metav{A}$ é necessariamente verdadeira, então não poderia ter \emph{deixado} de ser necessariamente verdadeira.

Isso nos leva a outro sistema novo, \mlSfour, que obtemos acrescentando a seguinte regra a~\mlT:
\factoidbox{
	\begin{fitchproof}
			\have[m]{m}{\ebox\metav{A}}
			\open
			\hypo[\ ]{k}{\ebox}\AS
			\have[n]{n}{\ebox\metav{A}}\rfour{m}
			\close
	\end{fitchproof}
Observe que R$\mathbf{4}$ só pode ser aplicada se a linha $m$ (que contém $\ebox A$) estiver fora da subprova estrita em que a linha $n$ ocorre e se essa subprova estrita não fizer parte de outra subprova estrita que não contenha~$m$.
}

A regra R$\mathbf{4}$ se parece com a regra {\ebox}E, exceto pelo fato de que, em vez de produzir $\metav{A}$ a partir de $\ebox\metav{A}$, produz $\ebox\metav{A}$ dentro de uma subprova estrita. A restrição, contudo, é a mesma: R$\mathbf{4}$ permite \textquotedblleft importar\textquotedblright{} $\ebox \metav{A}$ para uma subprova estrita, mas não para uma subprova estrita aninhada dentro de outra subprova estrita. Porém, se isso for necessário, uma aplicação adicional de R$\mathbf{4}$ produzirá o mesmo resultado.

Agora podemos provar ainda mais resultados. Por exemplo:
\begin{fitchproof}
	\open
	\hypo{1}{\ebox A}\AS
	\open
	\hypo{2}{\ebox}\AS
	\have{3}{\ebox A}\rfour{1}
	\close
	\have{4}{\ebox\ebox A}\boxi{2-3}
	\close
	\have{5}{\ebox A \eif \ebox\ebox A}\ci{1-6}
\end{fitchproof}
De modo semelhante, podemos provar $\ediamond\ediamond A \eif \ediamond A$. Isso mostra que, além de permitir \emph{acrescentar} caixas extras, \mlSfour{} permite \emph{eliminar} diamantes extras: de $\ediamond\ediamond \metav{A}$, você sempre pode inferir $\ediamond\metav{A}$.

\section{Sistema \mlSfive}
\label{S5}

Em \mlSfour, podemos sempre acrescentar uma caixa diante de outra caixa. Mas \mlSfour{} não nos permite automaticamente acrescentar uma caixa diante de um \emph{diamante}. Isto é, \mlSfour{} em geral não permite a inferência de $\ediamond\metav{A}$ para $\ebox\ediamond\metav{A}$. Mas, novamente, isso pode parecer uma deficiência, pelo menos se estivermos lendo $\ebox$ e $\ediamond$ como expressando \emph{necessidade} e \emph{possibilidade}. Parece intuitivo que, se $\metav{A}$ é possivelmente verdadeira, então não poderia ter \emph{deixado} de ser possivelmente verdadeira.

Isso nos leva ao nosso sistema modal final, \mlSfive, que obtemos acrescentando a seguinte regra à \mlSfour:
\factoidbox{
	\begin{fitchproof}
			\have[m]{m}{\enot \ebox\metav{A}}
			\open
			\hypo[\ ]{k}{\ebox}\AS
			\have[n]{n}{\enot\ebox\metav{A}}\rfive{m}
			\close
	\end{fitchproof}
A regra R$\mathbf{5}$ só pode ser aplicada se a linha $m$ (que contém $\enot \ebox\metav{A}$) estiver fora da subprova estrita em que a linha $n$ ocorre e se essa subprova estrita não fizer parte de outra subprova estrita que não contenha a linha~$m$.
}

Essa regra nos permite mostrar, por exemplo, que $\ediamond\ebox A\proves_\mlSfive  \ebox A:
\begin{fitchproof}
	\hypo{1}{\ediamond\ebox A}\PR
	\have{2}{\enot\ebox\enot\ebox A}\by{Def\ediamond}{1}
	\open
	\hypo{3}{\enot \ebox A}\AS
	\open
	\hypo{4}{\ebox}\AS
	\have{5}{\enot\ebox A}\rfive{3}
	\close
	\have{6}{\ebox\enot\ebox A}\boxi{4-5}
	\have{7}{\ered}\ne{2,6}
	\close
	\have{8}{\ebox A}\by{IP}{3-7}
\end{fitchproof}

Assim, além de acrescentar caixas diante de diamantes, também podemos eliminar diamantes diante de caixas.

Obtivemos \mlSfive{} acrescentando a regra~R$\mathbf{5}$ à \mlSfour. Na
verdade, poderíamos ter acrescentado a regra~R$\mathbf{5}$ à \mlT{}, deixado de fora a
regra~R$\mathbf{4}$ e obtido um sistema equivalente. Isso ocorre porque
tudo o que podemos provar usando a regra~R$\mathbf{4}$ também pode ser provado
usando R\mlT{} junto com~R$\mathbf{5}$. Por exemplo, eis uma
prova que mostra $\ebox A \proves_\mlSfive \ebox\ebox A$ sem
usar~R$\mathbf{4}:
\begin{fitchproof}
	\hypo{1}{\ebox A}\PR
	\open
	\hypo{2}{\ebox\enot\ebox A}\AS
	\have{3}{\enot\ebox A}\rt{2}
	\have{4}{\ered}\ne{1,3}
	\close
	\have{5}{\enot\ebox\enot\ebox A}\ni{2-4}
	\open
	\hypo{6}{\ebox}\AS
	\open
	\hypo{7}{\enot\ebox A}\AS
	\open
	\hypo{8}{\ebox}\AS
	\have{9}{\enot\ebox A}\rfive{7}
	\close
	\have{10}{\ebox \enot\ebox A}\boxi{8-9}
	\have{11}{\enot\ebox\enot\ebox A}\rfive{5}
	\have{12}{\ered}\ne{10,11}
	\close
	\have{13}{\ebox A}\ip{7-12}
	\close
	\have{14}{\ebox\ebox A}\boxi{6-13}
\end{fitchproof}
\mlSfive{} é \emph{estritamente mais forte} que \mlSfour: há coisas que podem ser provadas em \mlSfive, mas não em \mlSfour{} (por exemplo, $\ediamond\ebox A \eif \ebox A).

O ponto importante sobre \mlSfive{} pode ser expresso assim: se você tem uma longa sequência de caixas e diamantes, em qualquer combinação, pode eliminar todos menos o último. Por exemplo, $\ediamond\ebox\ediamond\ediamond\ebox\ebox\ediamond\ebox A$ pode ser simplificada para simplesmente $\ebox A.

\practiceproblems

\problempart
Forneça provas para os seguintes:
\begin{compactlist}
	\item $\ebox (A\eand B)\proves_\mlK\ebox A \eand \ebox B$
	\item $\ebox A\eand\ebox B\proves_\mlK\ebox( A \eand  B)$
	\item $\ebox A\eor\ebox B\proves_\mlK\ebox( A \eor  B)$
	\item $\ebox (A \eiff B)\proves_\mlK \ebox A \eiff \ebox B$
\end{compactlist}

\problempart
Forneça provas para os seguintes (sem usar conversão modal!):
\begin{compactlist}
	\item $\enot\ebox A\proves_\mlK \ediamond \enot A$
	\item $\ediamond\enot A\proves_\mlK\enot \ebox A$
	\item $\enot\ediamond A\proves_\mlK\ebox\enot A$
	\item $\ebox\enot A\proves_\mlK\enot\ediamond A$
\end{compactlist}

\problempart
Forneça provas para os seguintes (agora você pode usar conversão modal!):
\begin{compactlist}
	\item $\ebox(A\eif B), \ediamond A \proves_\mlK \ediamond B$
	\item $\ebox A \proves_\mlK \enot\ediamond\enot A$
	\item $\enot\ediamond\enot A \proves_\mlK \ebox A$
\end{compactlist}

\problempart
Forneça provas para os seguintes:
\begin{compactlist}
	\item $P\proves_\mlT\ediamond P$
	\item $\proves_\mlT (A\eand B)\eor(\enot \ebox A\eor\enot\ebox B)$
\end{compactlist}

\problempart
Forneça provas para os seguintes:
\begin{compactlist}
	\item $\ebox(\ebox A\eif B), \ebox (\ebox B\eif C), \ebox A \proves_\mlSfour \ebox\ebox C$
	\item $\ebox A \proves_\mlSfour \ebox(\ebox A \eor B)$
	\item $\ediamond \ediamond A \proves_\mlSfour \ediamond A$
\end{compactlist}


\problempart
Forneça provas em \mlSfive{} para os seguintes:
\begin{compactlist}
	\item $\enot\ebox\enot A, \ediamond B\proves_\mlSfive \ebox(\ediamond A \eand \ediamond B)$
	\item $A \proves_\mlSfive  \ebox\ediamond A$
	\item $\ediamond\ediamond A\proves_\mlSfive  \ediamond A$
\end{compactlist}


\chapter{Semantics for ML}
\label{Semantics}

So far, we have focussed on laying out various systems of Natural Deduction for ML. Now we will look at the \emph{semantics} for ML. A semantics for a language is a method for assigning truth-values to the sentences in that language. So a semantics for ML is a method for assigning truth-values to the sentences of ML.

\section{Interpretations of ML}

The big idea behind the semantics for ML is this. In ML, sentences are not just true or false, full stop. A sentence is true or false \emph{at a given possible world}, and a single sentence may well be true at some worlds and false at others. We then say that $\ebox \metav{A}$ is true \ifeff{} $\metav{A}$ is true at \emph{every} world, and $\ediamond\metav{A}$ is true \ifeff{} $\metav{A}$ is true at \emph{some} world.

That's the big idea, but we need to refine it and make it more precise. To do this, we need to introduce the idea of an \emph{interpretation} of ML. The first thing you need to include in an interpretation is a collection of \emph{possible worlds}. Now, at this point you might well want to ask: What exactly is a possible world? The intuitive idea is that a possible world is another way that this world could have been. But what exactly does that mean? This is an excellent philosophical question, and we will look at it in a lot of detail later. But we do not need to worry too much about it right now. As far as the formal logic goes, possible worlds can be anything you like. All that matters is that you supply each interpretation with a non-empty collection of things labelled \define{possible worlds}.

Once you have chosen your collection of possible worlds, you need to find some way of determining which sentences of ML are true at which possible worlds. To do that, we need to introduce the notion of a \emph{valuation function}. Those of you who have studied some maths will already be familiar with the general idea of a function. But for those of you who haven't, a function is a mathematical entity which maps arguments to values. That might sound a little bit abstract, but some familiar examples will help. Take the function $x+1$. This is a function which takes in a number as argument, and then spits out the next number as value. So if you feed in the number $1$ as an argument, the function $x+1$ will spit out the number $2$ as a value; if you feed in $2$, it will spit out $3$; if you feed in $3$, it will spit out $4$ \dots  Or here is another example: the function $x+y$. This time, you have to feed two arguments into this function if you want it to return a value: if you feed in $2$ and $3$ as your arguments, it spits out $5$; if you feed in $1003$ and $2005$, it spits out $3008$; and so on.

A valuation function for ML takes in a sentence and a world as its arguments, and then returns a truth-value as its value. So if $\nu$ is a valuation function and $w$ is a possible world, $\nu_w(\metav{A})$ is whatever truth-value $\nu$ maps $\metav{A}$ and $w$ to: if $\nu_w(\metav{A})=F$, then $\metav{A}$ is false at world $w$ on valuation $\nu$; if $\nu_w(\metav{A})=T$, then $\metav{A}$ is true at world $w$ on valuation $\nu$.

These valuation functions are allowed to map any \emph{atomic} sentence to any truth-value at any world. But there are rules about which truth-values more complex sentences get assigned at a world. Here are the rules for the connectives from TFL:
\begin{compactlist}
	\item $\nu_w(\enot\metav{A})=T$ \ifeff{} $\nu_w(\metav{A})=F$
	\item $\nu_w(\metav{A}\eand\metav{B})=T$ \ifeff{} $\nu_w(\metav{A})=T$ and $\nu_w(\metav{B})=T$
	\item $\nu_w(\metav{A}\eor\metav{B})=T$ \ifeff{} $\nu_w(\metav{A})=T$ or $\nu_w(\metav{B})=T$, or both
	\item $\nu_w(\metav{A}\eif\metav{B})=T$ \ifeff{} $\nu_w(\metav{A})=F$ or $\nu_w(\metav{B})=T$, or both
	\item $\nu_w(\metav{A}\eiff\metav{B})=T$ \ifeff{} $\nu_w(\metav{A})=T$ and $\nu_w(\metav{B})=T$, or $\nu_w(\metav{A})=F$ and $\nu_w(\metav{B})=F$
\end{compactlist}
So far, these rules should all look very familiar. Essentially, they all work exactly like the truth-tables for TFL. The only difference is that these truth-table rules have to be applied over and over again, to one world at a time.

But what are the rules for the new modal operators, $\ebox$ and $\ediamond$? The most obvious idea would be to give rules like these:
\begin{itemize}
	\item $\nu_w(\ebox \metav{A})=T$ \ifeff{} $\forall w' (\nu_{w'}(\metav{A})=T)$
	\item $\nu_w(\ediamond \metav{A})=T$ \ifeff{} $\exists w' (\nu_{w'}(\metav{A})=T)$
\end{itemize}
This is just the fancy formal way of writing out the idea that $\ebox\metav{A}$ is true at $w$ just in case $\metav{A}$ is true at \emph{every} world, and $\ediamond\metav{A}$ is true at $w$ just in case $\metav{A}$ is true at \emph{some} world.

However, while these rules are nice and simple, they turn out not to be quite as useful as we would like. As we mentioned, ML is meant to be a very flexible tool. It is meant to be a general framework for dealing with lots of different kinds of necessity. As a result, we want our semantic rules for $\ebox$ and $\ediamond$ to be a bit less rigid. We can do this by introducing another new idea: \emph{accessibility relations}.

An accessibility relation, $R$, is a relation between possible worlds. Roughly, to say that $Rw_1w_2$ (in English: world $w_1$ \emph{accesses} world $w_2$) is to say that $w_2$ is possible \emph{relative to} $w_1$. In other words, by introducing accessibility relations, we open up the idea that a given world might be possible relative to some worlds but not others. This turns out to be a \emph{very} fruitful idea when studying modal systems. We can now give the following semantic rules for $\ebox$ and $\ediamond$:
\begin{compactlist}[resume]
	\item $\nu_{w_1}(\ebox \metav{A})=T$ \ifeff{} $\forall w_2 (Rw_1w_2\eif \nu_{w_2}(\metav{A})=T)$
	\item $\nu_{w_1}(\ediamond \metav{A})=T$ \ifeff{} $\exists w_2 (Rw_1w_2\eand \nu_{w_2}(\metav{A})=T)$
\end{compactlist}
Or in plain English: $\ebox\metav{A}$ is true in world $w_1$ \ifeff{} $\metav{A}$ is true in every world that is possible relative to $w_1$; and $\ediamond\metav{A}$ is true in world $w_1$ \ifeff{} $\metav{A}$ is true in some world that is possible relative to $w_1$.

So, there we have it. An interpretation for ML consists of three things: a collection of possible worlds, $W$; an accessibility relation, $R$; and a valuation function, $\nu$. The collection of `possible worlds' can really be a collection of anything you like. It really doesn't matter, so long as $W$ isn't empty. (For many purposes, it is helpful just to take a collection of numbers to be your collection of worlds.) And for now, at least, $R$ can be any relation between the worlds in $W$ that you like. It could be a relation which every world in $W$ bears to every world in $W$, or one which no world bears to any world, or anything in between. And lastly, $\nu$ can map any atomic sentence of ML to any truth-value at any world. All that matters is that it follows the rules (1)--(7) when it comes to the more complex sentences.

Let's look at an example. It is often helpful to present interpretations of ML as diagrams, like this:
\begin{center}
	\begin{tikzpicture}[alt={The numbers 1 and 2 are connected by arrows. There are arrows from 1 to itself, and from 2 to 1. Node 1 is labelled by A and not-B. Node 2 is labelled by not-A and B.}]
		\node (atom1) at (0,1) {1};
		\node (atom2) at (3,1) {2};
		\node (atom3) at (-0.25,0) {$A$};
		\node (atom4) at (3,0) {$\enot A$};
		\node (atom5) at (-0.25,-0.7) {$\enot B$};
		\node (atom6) at (3,-0.7) {$B$};
		\draw[->, thick] (atom1)+(-0.15,0.15) arc (-330:-30:.3);
		%\draw[->, thick] (atom2)+(0.15,-0.15) arc (-150:150:.3); 
		\draw[<-, thick] (atom1) -- (atom2);
		\draw (-0.8,-1.2) rectangle (0.4,0.5);
		\draw (2.4,-1.2) rectangle (3.6,0.5);
	\end{tikzpicture}
	\bmlDescription{The numbers 1 and 2 are connected by arrows. There are arrows from 1 to itself, and from 2 to 1. Node 1 is labelled by A and not-B. Node 2 is labelled by not-A and B.}
\end{center}
Here is how to read the interpretation off from this diagram. It contains just two worlds, 1 and 2. The arrows between the worlds indicate the accessibility relation. So 1 and 2 both access 1, but neither 1 nor 2 accesses 2. The boxes at each world let us know which atomic sentences are true at each world: $A$ is true at 1 but false at 2; $B$ is false at 1 but true at 2. You may only write an atomic sentence or the negation of an atomic sentence into one of these boxes. We can figure out what truth-values the more complex sentences get at each world from that. For example, on this interpretation all of the following sentences are true at $w_1$:
\[
A\eand\enot B,\qquad B\eif A,\qquad \ediamond A,\qquad \ebox\enot B
\]
If you don't like thinking diagrammatically, then you can also present an interpretation like this:
\begin{ekey}
	\item[W]$1,2$
	\item[R]$\langle 1,1\rangle, \langle 2,1\rangle$
	\item[\nu]$\nu_{1}(A)=T, \nu_{1}(B)=F, \nu_{2}(A)=F, \nu_{2}(B)=T$
\end{ekey}
You will get the chance to cook up some interpretations of your own shortly, when we start looking at \emph{counter-interpretations}.

\section{A Semantics for System \mlK}
\label{SemanticsK}

We can now extend all of the semantic concepts of TFL to cover ML:
\factoidbox{
	\begin{factoiditems}
		\item $\metav{A}_1,\metav{A}_2, \dots
		\metav{A}_n\therefore\metav{C}$ is \define{modally valid} \ifeff{}
		there is no world in any interpretation at which
		$\metav{A}_1,\metav{A}_2, \dots \metav{A}_n$ are all true and
		$\metav{C}$ is false.

		\item $\metav{A}$ is a \define{modal truth} \ifeff{} $\metav{A}$
		is true at every world in every interpretation.

		\item $\metav{A}$ is a \define{modal contradiction} \ifeff{}
		$\metav{A}$ is false at every world in every interpretation.

		\item $\metav{A}$ is \define{modally satisfiable} \ifeff{}
		$\metav{A}$ is true at some world in some interpretation.
	\end{factoiditems}
		}
(From now on we will drop the explicit `modal' qualifications, since they can be taken as read.)

We can also extend our use of $\entails$. However, we need to add subscripts again, just as we did with $\proves$. So, when we want to say that $\metav{A}_1,\metav{A}_2, \dots \metav{A}_n\therefore\metav{C}$ is valid, we will write: $\metav{A}_1,\metav{A}_2, \dots \metav{A}_n\entails_\mlK\metav{C}$.

Let's get more of a feel for this semantics by presenting some counter-interpretations. Consider the following (false) claim:
\[\enot A\entails_\mlK \enot \ediamond A\]
In order to present a counter-interpretation to this claim, we need to cook up an interpretation which makes $\enot A$ true at some world $w$, and $\enot\ediamond A$ false at $w$. Here is one such interpretation, presented diagrammatically:
\begin{center}
	\begin{tikzpicture}[alt={The numbers 1 and 2 are connected by an arrow from 1
to 2. Node 1 is labelled not-A and node 2 is labelled A.}]
		\node (atom1) at (0,1) {1};
		\node (atom2) at (3,1) {2};
		\node (atom3) at (-0.25,0) {$\enot A$};
		\node (atom4) at (3,0) {$A$};
		\draw[->, thick] (atom1) -- (atom2);
		\draw (-0.8,-0.6) rectangle (0.4,0.5);
		\draw (2.4,-0.6) rectangle (3.6,0.5);
	\end{tikzpicture}
\bmlDescription{The numbers 1 and 2 are connected by an arrow from 1
to 2. Node 1 is labelled not-A and node 2 is labelled A.}
\end{center}
It is easy to see that this will work as a counter-interpretation for our claim. First, $\enot A$ is true at world $1$. And second, $\enot\ediamond A$ is false at $1$: $A$ is true at $2$, and $2$ is accessible from $1$. So there is some world in this interpretation where $\enot A$ is true and $\enot\ediamond A$ is false, so it is not the case that $\enot A\entails_\mlK\enot\ediamond A$.

Why did we choose the subscript \mlK? Well, it turns out that there is an important relationship between system \mlK{} and the definition of validity we have just given. In particular, we have the following two results:
\begin{compactlist}
	\item If $\metav{A}_1,\metav{A}_2, \dots \metav{A}_n\proves_\mlK\metav{C}$, then $\metav{A}_1,\metav{A}_2, \dots \metav{A}_n\entails_\mlK\metav{C}$
	\item If $\metav{A}_1,\metav{A}_2, \dots \metav{A}_n\entails_\mlK\metav{C}$, then $\metav{A}_1,\metav{A}_2, \dots \metav{A}_n\proves_\mlK\metav{C}$
\end{compactlist}
The first result is known as a \emph{soundness} result, since it tells us that the rules of \mlK{} are good, sound rules: if you can vindicate an argument by giving a proof for it using system \mlK, then that argument really is valid. The second result is known as a \emph{completeness} result, since it tells us that the rules of \mlK{} are broad enough to capture all of the valid arguments: if an argument is valid, then it will be possible to offer a proof in \mlK{} which vindicates it.

Now, it is one thing to state these results, quite another to prove them. However, we will not try to prove them here. But the idea behind the proof of soundness will perhaps make clearer how strict subproofs work.

In a strict subproof, we are not allowed to make use of any information from outside the strict subproof, except what we import into the strict subproof using \ebox E. If we've assumed or proved $\ebox \metav{A}$, by \ebox E, we can used $\metav{A}$ inside a strict subproof. And in \mlK, that is the only way to import a formula into a strict subproof. So everything that can be proved inside a strict subproof must follow from formulas $\metav{A}$ where outside the strict subproof we have $\ebox \metav{A}$. Let's imagine that we are reasoning about what's true in a possible world in some interpretation. If we know that $\ebox\metav{A}$ is true in that possible world, we know that $\metav{A}$ is true in all accessible worlds. So, everything proved inside a strict subproof is true in all accessible possible worlds. That is why \ebox I is a sound rule.

\section{A Semantics for System \mlT}
\label{SemanticsT}

A few moments ago, we said that system \mlK{} is sound and complete. Where does that leave the other modal systems we looked at, namely  \mlT, \mlSfour{} and \mlSfive? Well, they are all \emph{unsound}, relative to the definition of validity we gave above. For example, all of these systems allow us to infer $A$ from $\ebox A$, even though $\ebox A\nentails_\mlK A$.

Does that mean that these systems are a waste of time? Not at all! These systems are only unsound \emph{relative to the definition of validity we gave above}. (Or to use symbols, they are unsound relative to $\entails_\mlK$.) So when we are dealing with these stronger modal systems, we just need to modify our definition of validity to fit. This is where accessibility relations come in really handy.

When we introduced the idea of an accessibility relation, we said that it could be any relation between worlds that you like: you could have it relating every world to every world, no world to any world, or anything in between. That is how we were thinking of accessibility relations in our definition of $\entails_\mlK$. But if we wanted, we could start putting some restrictions on the accessibility relation. In particular, we might insist that it has to be \emph{reflexive}:
\[\forall w\, Rww\]
In English: every world accesses itself. Or in terms of relative possibility: every world is possible relative to itself. If we imposed this restriction, we could introduce a new consequence relation, $\entails_\mlT$, as follows:
\factoidbox{
	$\metav{A}_1,\metav{A}_2, \dots \metav{A}_n\entails_\mlT \metav{C}$ \ifeff{} there is no world in any interpretation \emph{which has a reflexive accessibility relation}, at which $\metav{A}_1,\metav{A}_2, \dots \metav{A}_n$ are all true and $\metav{C}$ is false
}
We have attached the \mlT{} subscript to $\entails$ because it turns out that system \mlT{} is sound and complete relative to this new definition of validity:
\begin{compactlist}
	\item If $\metav{A}_1,\metav{A}_2, \dots \metav{A}_n\proves_\mlT\metav{C}$, then $\metav{A}_1,\metav{A}_2, \dots \metav{A}_n\entails_\mlT\metav{C}$
	\item If $\metav{A}_1,\metav{A}_2, \dots \metav{A}_n\entails_\mlT\metav{C}$, then $\metav{A}_1,\metav{A}_2, \dots \metav{A}_n\proves_\mlT\metav{C}$
\end{compactlist}
As before, we will not try to prove these soundness and completeness results. However, it is relatively easy to see how insisting that the accessibility relation must be reflexive will vindicate the R\mlT{} rule:
\factoidbox{
	\begin{fitchproof}
			\have[m]{m}{\ebox \metav{A}}
			\have[\, ]{n}{\metav{A}}\rt{m}
	\end{fitchproof}
}
To see this, just imagine trying to cook up a counter-interpretation to this claim:
\[
	\ebox \metav{A}\entails_\mlT \metav{A}
\]
We would need to construct a world, $w$, at which $\ebox \metav{A}$ was true, but $\metav{A}$ was false. Now, if $\ebox \metav{A}$ is true at $w$, then $\metav{A}$ must be true at every world $w$ accesses. But since the accessibility relation is reflexive, $w$ accesses $w$. So $\metav{A}$ must be true at $w$. But now $\metav{A}$ must be true \emph{and} false at $w$. Contradiction!

\section{A Semantics for \mlSfour}
\label{SemanticsS4}

How else might we tweak our definition of validity? Well, we might also stipulate that the accessibility relation has to be \emph{transitive}:
\[\forall w_1\forall w_2\forall w_3 ((Rw_1w_2 \eand Rw_2w_3)\eif Rw_1w_3)\]
In English: if $w_1$ accesses $w_2$, and $w_2$ accesses $w_3$, then $w_1$ accesses $w_3$. Or in terms of relative possibility: if $w_3$ is possible relative to $w_2$, and $w_2$ is possible relative to $w_1$, then $w_3$ is possible relative to $w_1$. If we added this restriction on our accessibility relation, we could introduce a new consequence relation, $\entails_\mlSfour$, as follows:
\factoidbox{
	$\metav{A}_1,\metav{A}_2, \dots \metav{A}_n\entails_\mlSfour \metav{C}$ \ifeff{} there is no world in any interpretation \emph{which has a reflexive and transitive accessibility relation}, at which $\metav{A}_1,\metav{A}_2, \dots \metav{A}_n$ are all true and $\metav{C}$ is false
}
We have attached the \mlSfour{} subscript to $\entails$ because it turns out that system \mlSfour{} is sound and complete relative to this new definition of validity:
\begin{compactlist}
	\item If $\metav{A}_1,\metav{A}_2, \dots \metav{A}_n\proves_\mlSfour\metav{C}$, then $\metav{A}_1,\metav{A}_2, \dots \metav{A}_n\entails_\mlSfour\metav{C}$
	\item If $\metav{A}_1,\metav{A}_2, \dots \metav{A}_n\entails_\mlSfour\metav{C}$, then $\metav{A}_1,\metav{A}_2, \dots \metav{A}_n\proves_\mlSfour\metav{C}$
\end{compactlist}
As before, we will not try to prove these soundness and completeness results. However, it is relatively easy to see how insisting that the accessibility relation must be transitive will vindicate the \mlSfour{} rule:
\factoidbox{
	\begin{fitchproof}
			\have[m]{m}{\ebox\metav{A}}
			\open
			\hypo[\ ]{k}{\ebox}\AS
			\have[\ ]{n}{\ebox\metav{A}}\rfour{m}
	\end{fitchproof}
}
The idea behind strict subproofs, remember, is that they are ways to prove things that must be true in all accessible worlds. So the R$\mathbf{4}$ rule means that whenever $\ebox \metav{A}$ is true, $\ebox \metav{A}$ must also be true in every accessible world. In other words, we must have $\ebox\metav{A} \entails_\mlSfour \ebox\ebox\metav{A}$.

To see this, just imagine trying to cook up a counter-interpretation to this claim:
\[\ebox\metav{A} \entails_\mlSfour \ebox \ebox \metav{A}\]
We would need to construct a world, $w_1$, at which $\ebox\metav{A}$ was true, but $\ebox \ebox \metav{A}$ was false. Now, if $\ebox \ebox \metav{A}$ is false at $w_1$, then $w_1$ must access some world, $w_2$, at which $\ebox\metav{A}$ is false. Equally, if $\ebox \metav{A}$ is false at $w_2$, then $w_2$ must access some world, $w_3$, at which $\metav{A}$ is false. We just said that $w_1$ accesses $w_2$, and $w_2$ accesses $w_3$. So since we are now insisting that the accessibility relation be transitive, $w_1$ must access $w_3$. And as $\ebox\metav{A}$ is true at $w_1$, and $w_3$ is accessible from $w_1$, it follows that $\metav{A}$ must be true at $w_3$. So $\metav{A}$ is true \emph{and} false at $w_3$. Contradiction!

\section{A Semantics for \mlSfive}
\label{SemanticsS5}

Let's put one more restriction on the accessibility relation. This time, let's insist that it must also be \emph{symmetric}:
\[\forall w_1\forall w_2(Rw_1w_2 \eif Rw_2w_1)\]
In English: if $w_1$ accesses $w_2$, then $w_2$ accesses $w_1$. Or in terms of relative possibility: if $w_2$ is possible relative to $w_1$, then $w_1$ is possible relative to $w_2$. Logicians call a relation that is reflexive, symmetric, and transitive an \emph{equivalence} relation. We can now define a new consequence relation, $\entails_\mlSfive $, as follows:
\factoidbox{
	$\metav{A}_1,\metav{A}_2, \dots \metav{A}_n\entails_\mlSfive  \metav{C}$ \ifeff{} there is no world in any interpretation \emph{whose accessibility relation is an equivalence relation}, at which $\metav{A}_1,\metav{A}_2, \dots \metav{A}_n$ are all true and $\metav{C}$ is false
}
We have attached the \mlSfive{} subscript to $\entails$ because it turns out that system \mlSfive{} is sound and complete relative to this new definition of validity:
\begin{compactlist}
	\item If $\metav{A}_1,\metav{A}_2, \dots \metav{A}_n\proves_\mlSfive \metav{C}$, then $\metav{A}_1,\metav{A}_2, \dots \metav{A}_n\entails_\mlSfive \metav{C}$
	\item If $\metav{A}_1,\metav{A}_2, \dots \metav{A}_n\entails_\mlSfive \metav{C}$, then $\metav{A}_1,\metav{A}_2, \dots \metav{A}_n\proves_\mlSfive \metav{C}$
\end{compactlist}
As before, we will not try to prove these soundness and completeness results here. However, it is relatively easy to see how insisting that the accessibility relation must be an equivalence relation will vindicate the R$\mathbf{5}$ rule:
\factoidbox{
	\begin{fitchproof}
			\have[m]{m}{\enot\ebox \metav{A}}
			\open
			\hypo[\ ]{k}{\ebox}\AS
			\have[\, ]{n}{\enot\ebox\metav{A}}\rfive{m}
	\end{fitchproof}
}
The rule says that if $\metav{A}$ is not necessary, i.e., false in some accessible world, it is also not necessary in any accessible possible world, i.e., we have $\enot\ebox \metav{A} \proves_\mlSfive  \ebox\enot\ebox \metav{A}$.

To see this, just imagine trying to cook up a counter-interpretation to this claim:
\[
	\enot\ebox\metav{A} \entails_\mlSfive  \ebox \enot\ebox \metav{A}
\]
We would need to construct a world, $w_1$, at which $\enot\ebox\metav{A}$ was true, but $\ebox \enot\ebox \metav{A}$ was false.
Now, if $\enot\ebox\metav{A}$ is true at $w_1$, then $w_1$ must access some world, $w_2$, at which $\metav{A}$ is false. Equally, if $\ebox \enot\ebox \metav{A}$ is false at $w_1$, then $w_1$ must access some world, $w_3$, at which $\enot\ebox \metav{A}$ is false. Since we are now insisting that the accessibility relation is an equivalence relation, and hence symmetric, we can infer that $w_3$ accesses $w_1$. Thus, $w_3$ accesses $w_1$, and $w_1$ accesses $w_2$. Again, since we are now insisting that the accessibility relation is an equivalence relation, and hence transitive, we can infer that $w_3$ accesses $w_2$. But earlier we said that $\enot\ebox \metav{A}$ is false at $w_3$, which implies that $\metav{A}$ is true at every world which $w_3$ accesses. So $\metav{A}$ is true \emph{and} false at $w_2$. Contradiction!

In the definition of $\entails_\mlSfive $, we stipulated that the accessibility relation must be an equivalence relation. But it turns out that there is another way of getting a notion of validity fit for \mlSfive. Rather than stipulating that the accessibility relation be an equivalence relation, we can instead stipulate that it be a \emph{universal} relation:
\[\forall w_1\forall w_2\, Rw_1w_2\]
In English: every world accesses every world. Or in terms of relative possibility: every world is possible relative to every world. Using this restriction on the accessibility relation, we could have defined $\entails_\mlSfive $ like this:
\factoidbox{
	$\metav{A}_1,\metav{A}_2, \dots \metav{A}_n\entails_\mlSfive  \metav{C}$ \ifeff{} there is no world in any interpretation \emph{which has a universal accessibility relation}, at which $\metav{A}_1,\metav{A}_2, \dots \metav{A}_n$ are all true and $\metav{C}$ is false.
}
If we defined $\entails_\mlSfive $ like this, we would still get the same soundness and completeness results for \mlSfive. What does this tell us? Well, it means that if we are dealing with a notion of necessity according to which \emph{every} world is possible relative to \emph{every} world, then we should use \mlSfive. What is more, most philosophers assume that the notions of necessity that they are most concerned with, like \emph{logical necessity} and \emph{metaphysical necessity}, are of exactly this kind. So \mlSfive{} is the modal system that most philosophers use most of the time.

\practiceproblems

\problempart
Present counter-interpretations to the following false claims:
\begin{compactlist}
	\item $\enot P \entails_\mlK \enot\ediamond P$
	\item $\ebox(P \eor Q)\entails_\mlK \ebox P \eor \ebox Q$
	\item $\entails_\mlK \enot \ebox (A\eand \enot A)$
	\item $\ebox A\entails_\mlK A$
\end{compactlist}

\problempart
Present counter-interpretations to the following false claims:
\begin{compactlist}
	\item $\ediamond A\entails_\mlSfour \ebox\ediamond A$
	\item $\ediamond A, \ebox (\ediamond A \eif B)\entails_\mlSfour\ebox B$
\end{compactlist}

\problempart
Present counter-interpretations to the following false claims:
\begin{compactlist}
	\item $\ebox (M\eif O),\ediamond M\entails_\mlT O$
	\item $\ebox A\entails_\mlT \ebox \ebox A$
\end{compactlist}

\section*{Further reading}

Modal logic is a large subfield of logic. We have only scratched the surface. If you want to learn more about modal logic, here are some textbooks you might consult.

\begin{itemize}
	\item George E. Hughes and Max J. Cresswell, \textit{A New Introduction
	to Modal Logic}, Oxford: Routledge, 1996.
	\item Graham Priest, \textit{An Introduction to Non-Classical Logic},
	2nd ed., Cambridge: Cambridge University Press, 2008.
	\item James W. Garson, \textit{Modal Logic for Philosophers}, 2nd
	ed., Cambridge: Cambridge University Press, 2013.
\end{itemize}

None of these authors formulate their modal proof systems in quite the way we did, but the closest formulation is given by Garson.
